\bigskip

\chapter{Cronograma} \label{cha:cronograma}

As atividades necesárias para o desenvolvimento do trabalho foram elencadas
de forma a apresentar uma evolução gradual, com prazos sendo estimados que
permitam que sejam realizados, além do projeto em si, também sua documentação
de uso.

Com isso, é possível agrupar as etapas de acordo com seus objetivos dentro do
projeto, respeitando suas ordenações lógicas, com atividades de concepção e
construção sendo alocadas nos 1\textordmasculine\ e 2\textordmasculine\ 
semestres (\autoref{tab:crono1sem} e \autoref{tab:crono2sem}), respectivamente.

\tabela{Cronograma TCC I (1\textordmasculine\ semestre)}{
    \begin{tabular}{l|c|c|c|c}
        \label{tab:crono1sem}
        Etapa                       & Mar & Abr & Mai & Jun \\ \hline \hline
        Revisão Bibliográfica       &  X  &  X  &     &     \\ \hline
        Levantamento de Requisitos  &  X  &  X  &  X  &     \\ \hline
        Projeto da PCI              &     &  X  &  X  &  X  \\ \hline
        Relatório final             &     &     &     &  X
    \end{tabular}
}{O Autor (2025)}{}{}{}

\tabela{Cronograma TCC II (2\textordmasculine\ semestre)}{
    \begin{tabular}{l|c|c|c|c|c}
        \label{tab:crono2sem}
        Etapa                          & Jul & Ago & Set & Out & Nov \\ \hline \hline
        Construção da PCI              &  X  &  X  &     &     &     \\ \hline
        Desevolvimento do Firmware     &  X  &  X  &  X  &     &     \\ \hline
        Desenvolvimento do Software    &     &  X  &  X  &     &     \\ \hline
        Integração e Testes            &     &     &  X  &  X  &  X  \\ \hline
        Documentação e Relatório Final &     &     &     &  X  &  X
    \end{tabular}
}{O Autor (2025)}{}{}{}
